\section{模型准备与统一数值方法}
\subsection{输入数据与边界重构}
圆柱初始半径 $R_0=0.02$~m、长度 $L=0.25$~m，初始状态为
\begin{equation}T(r,0)=28~{}^\circ\mathrm C,\qquad C(r,0)=2.55~\mathrm{kg/kg}.\label{eq:initial}\end{equation}
附件1包含241个环境时刻，间隔60~s；附件2包含145个半径时刻，间隔1800~s。采用 PCHIP 保形插值\cite{fritsch1980}，环境在4~h后保持末值，半径仅在附件覆盖的72~h内使用。环境测量波动不强行单调化；半径按非增序列重构，调整数记录在数据审计文件中。

\subsection{固定域控制方程与边界}
固定半径下采用
\begin{align}
\rho(C)c_p(C)\frac{\partial T}{\partial t}&=\frac1r\frac\partial{\partial r}\left(rk(C)\frac{\partial T}{\partial r}\right),\label{eq:heat}\\
\frac{\partial C}{\partial t}&=\frac1r\frac\partial{\partial r}\left(rD(T,C)\frac{\partial C}{\partial r}\right).\label{eq:mass}
\end{align}
该有效扩散模型以均匀干物质参考权重定义含水率储存；题给经验密度只进入热方程，所涉及的近似已在模型假设中说明。
轴线为零通量，侧表面为
\begin{align}
\partial_rT(0,t)=\partial_rC(0,t)&=0,\label{eq:center}\\
-k\partial_rT(R,t)&=h_T(T_s-T_e),\label{eq:heatbc}\\
-D\partial_rC(R,t)&=h_m(C_s-C_e).\label{eq:massbc}
\end{align}
以向外通量为正，环境更热时热流指向内部，药材更湿时水分流向外部。

\subsection{题给物性与状态耦合}
附录2给出第一问的常热物性及非线性扩散系数：
\begin{equation}\rho=820,\quad c_p=2600,\quad k=0.36,\quad D=7\times10^{-9}e^{-0.89/C}.\label{eq:prop1}\end{equation}
第二、三问使用附录3：
\begin{align}
\rho&=650+128C,&c_p&=1450+2736\frac C{C+1},\nonumber\\
k&=0.21+0.38\frac C{C+1},&D&=2.4\times10^{-3}e^{-0.45/C}e^{-3850/(T+273.15)}.\label{eq:prop23}
\end{align}
第四问使用附录4：
\begin{align}
\rho&=760+90C,&c_p&=1850+2150\frac C{C+1},\nonumber\\
k&=0.12+0.20\frac C{C+1},&D&=4.2\times10^{-4}e^{-0.30/C}e^{-3850/(T+273.15)}.\label{eq:prop4}
\end{align}
含水率决定储热与导热物性，温度和含水率共同决定有效扩散率，构成状态耦合。

\subsection{守恒控制体与非线性界面通量}
内部求解使用\SolverNodes{}个径向节点，取表面加密映射 $\xi_i=1-[1-i/(N-1)]^2$，$i=0,\ldots,N-1$，并令 $r_i=R(t)\xi_i$；固定域取 $R(t)=R_0$。题目要求的21个物理位置是输出网格，不限制内部离散。各界面采用实际相邻节点距离 $r_{i+1}-r_i$ 计算传导系数 $G_f=A_f k_f/(r_{i+1}-r_i)$ 或 $A_fD_f/(r_{i+1}-r_i)$。界面在相邻节点中点，中心及表面控制体边界分别截在0和当前半径：
\begin{equation}V_i=\pi(r_{i+1/2}^2-r_{i-1/2}^2),\qquad A_f=2\pi r_f.\label{eq:cv}\end{equation}
以 $\phi$ 表示 $T$ 或 $C$，全隐式控制体方程为
\begin{equation}
\frac{s_iV_i}{\Delta t}(\phi_i^{n+1}-\phi_i^n)=G_E(\phi_{i+1}^{n+1}-\phi_i^{n+1})-G_W(\phi_i^{n+1}-\phi_{i-1}^{n+1}),\label{eq:fvm}
\end{equation}
热方程取 $s=\rho c_p$，含水率方程取 $s=1$。表面直接加入 $A_Rh(\phi_e-\phi_s)$。内部界面通量在相邻控制体中符号相反，因此离散求和能消去内部交换\cite{patankar1980}。

对 $k$ 用相邻节点平均；对变化剧烈的 $D(T,C)=b(T)e^{-a/C}$，使用 Kirchhoff 原函数
\begin{equation}F_a(C)=Ce^{-a/C}+a\operatorname{Ei}(-a/C),\qquad F_a'(C)=e^{-a/C},\label{eq:kirchhoff}\end{equation}
构造界面扩散率
\begin{equation}
D_f=b(T_f)\frac{F_a(C_{i+1})-F_a(C_i)}{C_{i+1}-C_i},\qquad T_f=\frac{T_i+T_{i+1}}2.\label{eq:faceD}
\end{equation}
当两浓度相近时使用中点极限。该格式对浓度方向积分，但仍近似界面温度和空间几何，不能称为原方程的精确离散。其作用是避免低含水率薄层中简单算术平均过度放大界面扩散率，效果由后文加密试验验证。

\subsection{非线性迭代、时间推进与事件}
每一步由旧解开始 Picard 迭代：由当前含水率求热物性及新温度，再更新界面扩散率并求含水率。当温度迭代差小于 $10^{-8}$~$^\circ$C、含水率迭代差小于 $10^{-10}$~kg/kg时结束，最多30次，失败则报错。
第一、二问时间步为1~s；第三、四问前3~h也为1~s，此后为10~s。步长在边界延拓切换、输出时刻处截断，避免跨越数据处理节点。第二、三问同网格、同物性、同前段步长，因此其共同时间点可直接核对。

定义 $M(t)=\max_i C_i(t)$，若 $M^n>0.15\ge M^{n+1}$，令
\begin{equation}
\lambda=\frac{M^n-0.15}{M^n-M^{n+1}},\qquad \widehat t_d=t_n+\lambda\Delta t.\label{eq:eventinterp}
\end{equation}
同时插值达标剖面。该量是离散数值解的事件估计，不能保证真实连续解的阈值恰落在这两个时间层之间；总误差还包括空间、时间和模型误差。提交文件另按题意逐秒或逐分钟采样，最终完整分钟覆盖估计事件。
